\begin{longtable}{TD}
\caption{Observation, signal, and feature terms.}
\label{tab:terms-observations-features}\\

\toprule
Term & Definition \\
\midrule
\endfirsthead

\toprule
Term & Definition \\
\midrule
\endhead

\bottomrule
\endfoot

\termrow{observation}{
Any measurement, event, trace, profile sample, counter, log entry, or
domain-specific signal exposed by a queue, worker runtime, operating system,
container runtime, dependency, tracing system, or domain component.
}

\termrow{raw observation}{
An unnormalized measurement or event before interpretation. A raw observation
does not by itself have stable model meaning.
}

\termrow{metric}{
A time-series, counter, gauge, or histogram observation. Metrics are one kind
of observation; traces, logs, profiles, and domain events are observations too.
}

\termrow{counter}{
A metric that counts events, operations, bytes, retries, failures, or other
monotonic quantities. A counter may represent a symptom or outcome rather than
a root influence.
}

\termrow{gauge}{
A metric representing a current value, such as queue depth, active workers,
available connections, or current memory usage.
}

\termrow{histogram}{
A metric representing a distribution of observed values, such as durations,
latencies, sizes, or per-attempt costs.
}

\termrow{trace}{
An execution trace composed of spans. A trace exposes execution structure and
timing, but it does not by itself define task-kind-specific relevance,
causality, feature mapping, or profiling policy.
}

\termrow{span}{
A timed unit inside a trace representing an operation, processing stage,
dependency call, or instrumented region. A span is not the same as a task
attempt.
}

\termrow{log}{
A structured or unstructured event record emitted by a system. Logs may expose
status, error classes, domain events, or diagnostic information.
}

\termrow{runtime profile}{
A CPU, heap, allocation, lock, blocking, or runtime-level profile. A runtime
profile is distinct from a behavior profile.
}

\termrow{signal}{
An interpreted, aggregated, or derived observation that indicates a condition
relevant to the model. A signal is an intermediate layer between raw
observation and canonical feature.
}

\termrow{canonical feature}{
A typed and normalized model input derived from task metadata, execution
context, raw observations, or signals. Task Behavior Graphs operate over
canonical features rather than arbitrary raw metrics.
}

\termrow{semantic feature}{
An informal synonym for a feature with explicit semantic meaning. This paper
uses the term canonical feature when referring to model inputs with defined
type, range, unit, influence domain, and intended meaning.
}

\termrow{feature mapping}{
The process of converting raw observations, signals, task metadata, and context
into canonical features. Feature mapping is separate from graph evaluation.
}

\termrow{normalization}{
The process of converting values into a defined type, unit, scale, bucket, or
range. For example, network measurements may be normalized into an instability
score in $[0,1]$.
}

\termrow{aggregation}{
The process of combining observations over a time window, group, task kind,
attempt set, or distribution. Examples include sum, average, maximum, p95, and
p99.
}

\termrow{feature registry}{
A registry of canonical features, including names, types, ranges, units,
influence domains, and intended meaning.
}

\termrow{metric descriptor}{
A descriptor of a raw metric, including source, unit, label set, aggregation,
cardinality, and mapping to canonical features.
}

\termrow{feature descriptor}{
A descriptor of a canonical feature, including type, range, unit, semantic
meaning, influence domain, and expected use in the model.
}

\termrow{cardinality}{
The number of distinct time series or label combinations produced by a metric.
Cardinality affects telemetry overhead and storage cost.
}

\termrow{high-cardinality label}{
A metric label that can create a large or unbounded number of series, such as
task id, user id, object id, URL, or unbounded tenant label.
}

\termrow{observation window}{
The time window over which observations are aggregated or interpreted. For
example, network instability may be computed over the last 30 seconds.
}

\termrow{freshness}{
How recent and applicable an observation is to the current task attempt. Stale
observations may reduce confidence.
}

\end{longtable}
